% Chapter 1

\chapter{Introduction}
\label{Chapter1}
\lhead{Chapter 1. \emph{Introduction}}

\section{Background}
Large Language Model-based autonomous agents are being widely applied in a variety of domains for long-range tasks including software support, data management, documentation processing and process automation. In many cases these agents use static registries of pre-defined tools provided to them from their developer. While this approach has advantages relative to auditing and permissioning; it severely limits the ability of the agent to perform tasks that require intermediate parsers, converters, benchmarks and integrations that were not provided prior to the beginning of the task.

The goal of dynamic tool synthesis is to provide agents the ability to dynamically create new capabilities by executing small programs, examining their results, and using those results as a capability for performing a particular type of task. However, the generation of new capabilities changes the underlying security model. It is possible for generated capabilities to include maliciously intended destructive operations, command injection prompts, resource consumption attacks, or deliberately designed host escapes. Therefore, simply sandboxing the process hosting the agent is often insufficient protection against maliciously generated code created during an open-ended conversation.

IronClaw addresses the challenge of integrating dynamic tool synthesis with Firecracker microVM isolation. The host process manages policy, LLM context, and microVM lifecycle decisions. Untrusted tool code runs in a guest runtime within a microVM. Data exchange occurs over a controlled transport boundary, allowing the system to return tool results without providing generated code with direct access to host resources.

\section{Use Cases}
IronClaw is designed for scenarios where autonomous agents need to solve tasks
that go beyond what a fixed registry of tools can handle while keeping host
systems secure. Here are some examples of tradeoffs that motivate the design.

\subsection{Ad Hoc Data Transformation}
Agents frequently encounter files or streams in formats for which no
pre-installed tool exists. For example, users might ask the agent to extract
metrics from custom log formats, convert between domain specific serialization
formats, or aggregate columns from several inconsistent CSV files. IronClaw can
synthesize small parsers or transformers and run them inside micro VMs.
Boundaries of this VM ensure that malformed or malicious transformations cannot
read host filesystem or leak data.

\subsection{Sandboxed Analysis of Untrusted Artifacts}
Security and incident response workflows often require analyzing files that
themselves might be hostile such as suspicious scripts, downloaded executables
or email attachments Rather than exposing hosts to artifacts, agents can
generate analysis tools and send approved explicit inputs to guests and observe
bounded outputs. If generated code tries to breach containment, it stays
confined to disposable guest environment.

\subsection{Dynamic Integration of APIs and Services}
An agent may encounter an external API for which no purpose-built capability is
available. It can then generate a minimal HTTP client or SDK wrapper on demand.
IronClaw runs generated clients in the guest with network policies enforced by
the host and thus integrates without giving tools persistent network privileges
on the host.

\subsection{Iterative Code Generation and Testing}
Often software engineering tasks require writing, running and refining small
programs. IronClaw lets agents generate candidate code and execute unit tests or
benchmarks in the guest and observe failures and then iterate. Each execution
runs in an isolated guest so buggy code or hostile code cannot corrupt workspace
on the host or persist between iterations.

\section{Problem Statement}
This dissertation asks whether autonomous agents gain measurable utility by
synthesizing and reusing their own tools while maintaining strong separation
from the host. When no purpose-built capability is available, repeatedly
authoring task logic inline may consume additional model steps and tokens.
Dynamic synthesis may avoid that repeated effort, but generated code increases
the risk to the host. IronClaw makes this utility/security trade-off measurable
by comparing no synthesis with dynamic synthesis for utility, and ordinary
process execution with microVM execution for containment and systems cost.

\section{Objectives}
The objectives of this thesis are:
\begin{itemize}
    \item To design and implement a dynamic execution path for agent-generated
    tools in IronClaw.
    \item To isolate untrusted generated code using Firecracker microVMs and a
    strict host/guest boundary.
    \item To compare agents with dynamic tool synthesis against otherwise
    equivalent agents that must author task logic inline on every request.
    \item To evaluate containment against simulated malicious generated tools.
    \item To profile the latency and resource overhead introduced by microVM
    isolation.
\end{itemize}

\section{Research Questions}
This thesis is guided by the following research questions:
\begin{itemize}
    \item \textbf{RQ1: Utility.} On repeated, complex tasks for which no
    purpose-built tool is available, does dynamic tool synthesis reduce agent
    interaction steps and token consumption relative to authoring the required
    logic inline, without reducing answer correctness?
    \item \textbf{RQ2: Security and Systems.} What is the measurable security
    containment capability and performance overhead of Firecracker microVM
    isolation compared to traditional process-level sandboxing when executing
    agent-generated code?
\end{itemize}

\section{Hypotheses}
The thesis evaluates two hypotheses:
\begin{itemize}
    \item \textbf{H1: Utility.} On repeated tasks for which no purpose-built
    tool is available, policy-guided dynamic tool synthesis reduces model
    interaction steps and token consumption relative to generic capabilities
    alone, while preserving task correctness. The benefit increases with reuse.
    \item \textbf{H2: Security and Performance.} IronClaw's microVM architecture
    contains simulated host-breakout attempts from malicious synthesized tools,
    at the cost of measurable but acceptable initialization latency and memory
    overhead.
\end{itemize}
